%% MAIN

The rule baselines separate sharply, and the ordering is informative.  The
Fixed formation, which never moves, covers $@@FIXED_CK@@$ of the sensor-steps:
it is the floor that the initial deployment alone provides, and it is also
the controller that never satisfies both invariants at once.  The Greedy
recharge controller, which returns to a pad when its battery falls below a
threshold and otherwise holds station, reaches $@@GREEDY_CK@@$ by trading
coverage for energy.  The Reactive heuristic, which moves the swarm towards
the uncovered sensors, reaches $@@REACT_CK@@$ with full connectivity and a
service continuity of $@@REACT_SC@@$, and it is the strongest controller in
the comparison.

It is also the controller with the most information, and the comparison
should be read with that in mind.  The reactive heuristic is given the
\emph{exact} set of under-covered sensors, which corresponds to an unbounded
broadcast of the coverage state; the learned policy receives a fixed-size
summary of it.  Part of the gap between them is therefore a difference of
information rather than of control, and the resolution of the broadcast is a
parameter of the framework rather than a free choice.

The learned controller does not reach the heuristic, and the two learned
baselines locate the reason.  The Masked PPO with a local critic differs
from the proposed method in exactly one respect: each agent has its own
actor and its own critic on its own observation, instead of a
parameter-shared actor with a centralised critic, so the advantage it
receives is its own rather than a broadcast team signal.  The Independent
DQN arm is the value-based point of comparison.  Vanilla MAPPO, which
removes the framework's observation terms, reward shaping, charging
arbitration, reserve mechanism and action mask, shows how much of the
learned method's behaviour survives without them.

The last row is the deterministic version of the \emph{same} learned actor.
The gap between it and the sampled policy is the symmetry collapse
discussed in the next subsection: with a shared actor the argmax action is
the same for every UAV, so the deterministic policy translates the swarm as
a rigid body.  Eq.~\eqref{eq:policy} specifies a sampled policy, and the
sampled policy is the one reported as the proposed method; the argmax row is
reported because the difference is informative, not because it is the
better controller.

%% OBS

The first row is the observation of Eq.~\eqref{eq:obs} as originally
written.  Its most striking property is not in the coverage column but in
the ``Distinct'' column: the deterministic policy selects almost the same
action for every one of the eight UAVs, and the swarm's spread barely grows
beyond the radius of the initial cluster.  This is the shared-actor symmetry
collapse, and it is precisely the failure that Foerster et al.'s
condition~(ii) is about.  With an observation that does not identify the
agent, a parameter-shared network has no basis on which to specialise, and
the consequence is measurable rather than theoretical.

Adding the identity does what the condition predicts.  The number of
distinct actions rises, the swarm disperses, and the coverage rises with it
--- by $@@B_MINUS_A@@$, which is $@@B_SHARE@@$ of the distance between the
local-only observation and the reactive heuristic.  The condition is
therefore not a formality: it is worth a measurable amount of persistent
coverage, and a reader who omits it obtains a worse controller without
necessarily knowing why.

@@C_SENTENCE@@

None of this closes the gap to the reactive heuristic, and the reason is
instructive.  Dispersing is not the same as dividing the work: an actor
that has learned to send its eight instances in eight different directions
has not learned to send them to eight different sensors.  The policy is
better at finding sensors than at dividing the work among them, and the
credit assignment, not the observation, is what limits it.

%% ENERGY

Arm~B removes the energy constraint entirely: the pack is ten times the
nominal size, so no UAV can strand and none ever needs to return.  The
coverage rises, and the service continuity rises with it, but by far less
than the gap to the reactive heuristic.  The energy constraint is a real
cost, worth $@@ENERGY_DELTA@@$ of persistent coverage, and it is not the
explanation.  Arm~C removes the threshold-violation term of
Table~\ref{tab:reward}, and it makes almost no difference: the stranding
persists and the coverage is unchanged.

Both results point the same way.  The learned policy's failure is not that
it mismanages energy; it is that it does not know what to do.  A policy
that does not know what to do wanders, and a wandering swarm empties its
packs, so the stranding is a \emph{symptom} of the credit-assignment
limitation and not its cause.  Removing the constraint removes the symptom
and leaves the coverage where it was.

This is the point at which the framework earns its keep.  A comparison that
reported only coverage numbers would have had to choose between two
explanations --- a badly specified energy reward, or something more
fundamental --- and would have had no way to distinguish them, because the
symptom and the cause have the same ordering.  The controlled arms
distinguish them, and the conclusion is the one from which the remedy
follows: the controller needs a per-agent signal, not a better team signal
and not a better observation.

%% ABL

Every variant is trained from the same seed on the same configuration
except for the one component that is removed, so a difference is
attributable to that component.  The table is reported in full, including
the variants that make no measurable difference, because the claim this
paper supports is a claim about which components are load-bearing.  @@ABL_TEXT@@

%% SCEN

Each row is the mean over the variants of that scenario, with the proposed
controller and the reactive heuristic evaluated on the same scenes.  The
matrix is where the framework's value is most visible, because it shows
which regime a controller is good in rather than in which regime it was
tuned.  @@SCEN_TEXT@@
